% ============================================================================
% sec3_structural_model.tex — V7: Formalized structural model (RAND standard)
% Paper: Frequent Losers as Cover Bidders in Public Procurement
% Target: RAND Journal of Economics
% ============================================================================

\section{A Framework for Cover-Bidder Deployment}
\label{sec:structural}

We develop a framework for cover bidding in procurement auctions
that formalizes the cover-bidder deployment decision, characterizes
equilibrium cover-bid distributions, and generates testable predictions
that map directly to our empirical strategy
(Section~\ref{sec:empirical}). The framework has three aims: (i)~to
provide economic foundations for the frequent-loser screen; (ii)~to
specify a mixture likelihood that separates cover bids from genuine bids
in the data; and (iii)~to derive the market-selection prediction---that
cover bidding concentrates in competitive markets---as an endogenous
outcome of the cartel's optimization.

% ── 3.1 Environment ──────────────────────────────────────────────

\subsection{Environment}
\label{sec:environment}

Consider a procurement auction for item $i$ at purchasing unit $k$ in
year $t$. The procuring agency posts a reference price $r > 0$ and
selects the lowest bidder. There are $n \geq 1$ \emph{genuine} bidders,
each with private cost $c_j \sim F(\cdot \mid \mathbf{x}_j)$, where
$\mathbf{x}_j$ is a vector of observable cost shifters (firm size, CNAE
sector, geographic distance to PBU~$k$). Genuine bidders submit bids
$b_j = c_j + \mu_j$, where $\mu_j \geq 0$ is a markup that depends on
the competitive environment.

A cartel controls firm $j^*$---the designated winner---and deploys
$m \geq 0$ \emph{cover bidders} (frequent losers). Each cover bidder
submits a bid $b_\ell$ drawn from a distribution $G(\cdot \mid b^*,
\boldsymbol{\theta})$, where $b^*$ is the cartel's designated winning
bid and $\boldsymbol{\theta}$ is a vector of strategic parameters. Cover
bids are designed to \emph{not} win: $b_\ell > b^*$ with probability
one. The procuring agency observes $n + m$ total bids.

Two institutional features of S\~ao Paulo's BEC platform shape the
model:

\begin{enumerate}
  \item \textbf{Convite} (sealed-bid): requires at least three bidders
    (Lei~8.666/93, Art.~22, \S7). If $n < 3$, the cartel must deploy
    $m \geq 3 - n$ cover bidders to satisfy the minimum-bidder
    constraint. Bids are submitted simultaneously and opened together.
  \item \textbf{Preg\~ao} (electronic reverse auction): no
    minimum-bidder requirement. Bidders submit descending bids in real
    time. Cover bidders can observe the auction dynamics before
    calibrating their bids, facilitating coordination.
\end{enumerate}

\noindent We denote the auction modality by $\tau \in \{\text{convite},
\text{preg\~ao}\}$, the minimum-bidder constraint by $\underline{n}(\tau)$
(equal to 3 for convite, 0 for preg\~ao), and the detection probability
by $\theta_k \in [0,1]$, which varies across PBUs and is increasing in
oversight capacity.

% ── 3.2 Cartel's Problem ─────────────────────────────────────────

\subsection{The Cartel's Optimization Problem}
\label{sec:cartel_problem}

The cartel chooses the number of cover bidders $m$, the winning bid
$b^*$, and the cover-bid distribution $G$ to maximize expected profit
net of cover-bidding costs and detection penalties:
\begin{equation}
  \max_{m \in \mathbb{Z}_+,\, b^*,\, G} \; \underbrace{\pi(b^*, m, n)
  }_{\text{cartel surplus}}
  - \underbrace{C(m)}_{\text{cover cost}}
  - \underbrace{\theta_k \cdot \Phi(m)}_{\text{expected penalty}}
  \label{eq:cartel_objective}
\end{equation}
subject to:
\begin{align}
  b^* &\leq r \label{eq:ref_price_constraint} \\
  \Pr(b_\ell > b^* \mid G) &= 1 \quad \forall \, \ell \in \{1, \ldots, m\}
    \label{eq:cover_above} \\
  n + m &\geq \underline{n}(\tau) \label{eq:min_bidder}
\end{align}

\noindent We maintain the following assumptions throughout: % V7: explicit assumptions

\begin{description}
  \item[A1 (Surplus structure).] The cartel's gross surplus
    $\pi(b^*, m, n) = b^* - c_{j^*}$ is the gap between the winning bid
    and the designated winner's cost. We assume $\pi$ is weakly
    increasing in $m$: additional cover bidders help sustain $b^*$ by
    creating the appearance of competition and satisfying procurement
    regulations. Formally, $\partial \pi / \partial m \geq 0$.

  \item[A2 (Diminishing returns).] The marginal return to cover bidding
    is decreasing: $\partial^2 \pi / \partial m^2 < 0$. The first cover
    bidder has high value (satisfying the minimum-bidder rule or creating
    basic noise); subsequent cover bidders have diminishing value as the
    appearance of competition saturates.

  \item[A3 (Cover-bidding cost).] $C(m) = c_0 + c_1 m$ with $c_0 \geq 0$
    and $c_1 > 0$. The cost is linear in the number of cover bidders:
    $c_0$ is the fixed setup cost (e.g., establishing the coordination
    mechanism) and $c_1$ is the per-bidder cost (recruitment,
    compensation, coordination).

  \item[A4 (Detection penalty).] The expected penalty is
    $\theta_k \cdot \Phi(m)$ where $\Phi(m)$ is increasing and convex:
    $\Phi'(m) > 0$ and $\Phi''(m) \geq 0$. More cover bidders increase
    the forensic footprint and the probability of detection conditional
    on investigation. For the closed-form solution in
    Proposition~\ref{prop:optimal_m}, we use the linear specialization
    $\Phi(m) = \phi_0 m$ with $\phi_0 > 0$.%
    \footnote{The more general specification $\Phi(m, b^*) =
    \phi_0 m + \phi_1 (b^* - c_{j^*})$ nests markup-dependent penalties.
    Under this specification, the FOC for $b^*$ yields
    $1 - \theta_k \phi_1 = 0$, implying $b^* = c_{j^*} + 1/\phi_1$
    when $\theta_k \phi_1 < 1$ (interior solution) or $b^* = r$ when
    the reference-price constraint binds. This does not affect the
    optimal $m^*$ conditional on $b^*$, so the comparative statics in
    Proposition~\ref{prop:optimal_m} hold under the more general
    penalty.}

  \item[A5 (Interior solution exists).] $\partial \pi / \partial m
    \big|_{m=0} > c_1 + \theta_k \phi_0$, so that at least one cover
    bidder is profitable in the absence of the minimum-bidder constraint.
\end{description}

\begin{proposition}[Optimal number of cover bidders]
\label{prop:optimal_m}
Under Assumptions A1--A5, the cartel's optimal cover-bidder count is:
\begin{equation}
  m^* = \max\!\left\{\underline{n}(\tau) - n, \;
  \left\lfloor \frac{\partial \pi / \partial m}{c_1 + \theta_k \phi_0}
  \right\rfloor \right\}
  \label{eq:optimal_m}
\end{equation}
\noindent where the first argument is the corner solution from the
minimum-bidder constraint~\eqref{eq:min_bidder} and the second is the
interior optimum. The optimal $m^*$ satisfies:
\begin{enumerate}
  \item[(i)] $\partial m^* / \partial \theta_k < 0$\,: higher detection
    probability reduces cover bidding;
  \item[(ii)] $\partial m^* / \partial c_1 < 0$\,: higher per-bidder
    cost reduces cover bidding;
  \item[(iii)] $m^*$ is weakly larger under convite than preg\~ao,
    holding $n$ fixed, because $\underline{n}(\text{convite}) = 3 >
    0 = \underline{n}(\text{preg\~ao})$.
\end{enumerate}
\end{proposition}

\begin{proof}
Relaxing the integer constraint, the cartel's first-order condition
for $m$ is:
\begin{equation}
  \frac{\partial \pi}{\partial m} = c_1 + \theta_k \phi_0
  \label{eq:foc_m}
\end{equation}
By A2, $\partial^2 \pi / \partial m^2 < 0$, so the second-order
condition $\partial^2 \pi / \partial m^2 - \theta_k \Phi''(m) < 0$
is satisfied (both terms are non-positive). The interior solution
$\tilde{m}$ defined by~\eqref{eq:foc_m} is therefore unique. When
the constraint $n + m \geq \underline{n}(\tau)$ binds (i.e.,
$\underline{n}(\tau) - n > \tilde{m}$), the optimal solution is
the corner $m^* = \underline{n}(\tau) - n$. Otherwise,
$m^* = \lfloor \tilde{m} \rfloor$ (rounding to the nearest integer
below $\tilde{m}$). The stated formula follows.

For the comparative statics, differentiate~\eqref{eq:foc_m} implicitly.
Let $H(m, \theta_k) \equiv \partial \pi / \partial m - c_1 -
\theta_k \phi_0 = 0$ define the interior solution. Then:
\[
  \frac{\partial m^*}{\partial \theta_k} =
  -\frac{\partial H / \partial \theta_k}{\partial H / \partial m}
  = -\frac{-\phi_0}{\partial^2 \pi / \partial m^2}
  = \frac{-\phi_0}{\partial^2 \pi / \partial m^2} < 0
\]
since $\phi_0 > 0$ and $\partial^2 \pi / \partial m^2 < 0$ (A2).
The argument for $\partial m^* / \partial c_1 < 0$ is identical,
replacing $\phi_0$ with $1$. Property~(iii) follows from the
definition of $m^*$ as the maximum of the corner and interior
solutions: convite's higher $\underline{n}$ weakly raises the corner
solution.
\end{proof}

% ── 3.3 Cover-Bid Distributions ──────────────────────────────────

\subsection{Equilibrium Cover-Bid Distributions}
\label{sec:cover_distributions}

The cartel chooses the cover-bid distribution $G$ to minimize the
probability of detection while satisfying
constraint~\eqref{eq:cover_above}. We characterize two equilibrium
regimes that differ in the cartel's information and coordination
capacity.

\begin{proposition}[Regime~1: Complementary cover bidding]
\label{prop:regime1}
Suppose the cartel can instruct cover bidders to ``bid above $b^*$''
with support $[b^*, b^* + \delta]$ for some $\delta > 0$, but cannot
specify exact bid levels (limited coordination). Suppose further that
the detection algorithm is a variance screen that flags tenders where
the bid distribution departs from maximum entropy over the feasible
support.\footnote{A variance screen computes the Kullback--Leibler
divergence $D_{\text{KL}}(G \| U)$ between the submitted bid
distribution $G$ and the uniform distribution $U$ over $[b^*, b^* +
\delta]$. The cartel minimizes $D_{\text{KL}}$ subject to
$\text{supp}(G) \subseteq [b^*, b^* + \delta]$; the solution is
$G = U$.} Then the equilibrium cover-bid distribution is:
\begin{equation}
  G_1(b) = \frac{b - b^*}{\delta}, \quad b \in [b^*, b^* + \delta]
  \label{eq:regime1}
\end{equation}
Under Regime~1, $\text{Var}(b_\ell) = \delta^2 / 12$. If $\delta$
is large, $\sigma_{\text{cover}} > \sigma_{\text{genuine}}$.
\end{proposition}

\begin{proof}
The cartel minimizes the KL divergence between $G$ and the maximum-entropy
distribution over $[b^*, b^* + \delta]$. By Gibbs' inequality,
$D_{\text{KL}}(G \| U) \geq 0$ with equality if and only if $G = U$
almost everywhere. Since the constraint set
$\{G : \text{supp}(G) \subseteq [b^*, b^* + \delta]\}$ includes $U$,
the minimum is attained at $G_1 = U[b^*, b^* + \delta]$.
\end{proof}

\begin{proposition}[Regime~2: Coordinated cover bidding]
\label{prop:regime2}
Suppose the cartel observes $b^*$ in real time (feasible under
preg\~ao) and can specify cover bids precisely (full coordination).
Suppose the detection algorithm is a dispersion screen that flags
tenders where the within-tender bid standard deviation is anomalous.
Then the equilibrium cover-bid distribution solves:
\begin{equation}
  \min_{\sigma_c > 0} \; \left| \sigma_c - \sigma_{\text{genuine}}
  \right| \quad \text{s.t.} \quad \Pr(b_\ell > b^*) = 1
  \label{eq:regime2_opt}
\end{equation}
The solution is a truncated normal distribution centered at
$\mu_c = b^* + \epsilon$ with $\epsilon > 0$ and
$\sigma_c < \sigma_{\text{genuine}}$:
\begin{equation}
  G_2(b) = \Phi\!\left(\frac{b - \mu_c}{\sigma_c}\right),
  \quad \mu_c = b^* + \epsilon, \quad
  \sigma_c < \sigma_{\text{genuine}}
  \label{eq:regime2}
\end{equation}
\end{proposition}

\begin{proof}
The cartel seeks $\sigma_c$ as close to $\sigma_{\text{genuine}}$ as
possible (to mimic genuine bids) while ensuring $\Pr(b_\ell > b^*) = 1$.
For a normal distribution with mean $\mu_c$ and standard deviation
$\sigma_c$, the constraint $\Pr(b_\ell > b^*) = 1$ requires
$\mu_c - b^* \gg \sigma_c$ (the entire distribution lies above $b^*$).
In practice, this is satisfied by setting $\mu_c = b^* + \epsilon$
with $\epsilon > 0$ and truncating at $b^*$. The truncation compresses
the effective standard deviation of the observed bids below $\sigma_c$,
so the dispersion-minimizing solution has $\sigma_c$ set to match
$\sigma_{\text{genuine}}$ from below. The constraint $\mu_c > b^*$
(cover bids must exceed the winning price) prevents exact matching,
yielding $\sigma_c < \sigma_{\text{genuine}}$.
\end{proof}

% ── 3.4 Market Selection ─────────────────────────────────────────

\subsection{Market Selection: Why Cover Bidding Concentrates in
Competitive Markets}
\label{sec:market_selection}

\begin{proposition}[Market-selection prediction]
\label{prop:market_selection}
Under Assumptions A1--A2, if the marginal return to cover bidding is
decreasing in market concentration---formally, if
\begin{equation}
  \frac{\partial^2 \pi}{\partial m \, \partial \text{HHI}} < 0
  \label{eq:cross_deriv}
\end{equation}
where $\text{HHI}$ is the Herfindahl--Hirschman Index of genuine
competition---then $m^*$ is decreasing in HHI. Cover bidding
concentrates in competitive markets.
\end{proposition}

\begin{proof}
From the FOC~\eqref{eq:foc_m} at the interior solution:
$\partial \pi / \partial m = c_1 + \theta_k \phi_0$. The right-hand
side does not depend on HHI. Differentiating implicitly:
\[
  \frac{\partial^2 \pi}{\partial m^2} \cdot
  \frac{\partial m^*}{\partial \text{HHI}}
  + \frac{\partial^2 \pi}{\partial m \, \partial \text{HHI}} = 0
\]
Solving:
\[
  \frac{\partial m^*}{\partial \text{HHI}} =
  -\frac{\partial^2 \pi / \partial m \, \partial \text{HHI}}
  {\partial^2 \pi / \partial m^2}
\]
By A2, $\partial^2 \pi / \partial m^2 < 0$ (denominator negative).
By assumption~\eqref{eq:cross_deriv}, the numerator is negative.
Therefore $\partial m^* / \partial \text{HHI} < 0$.

The economic content of~\eqref{eq:cross_deriv} is that cover bidders
are more valuable in competitive markets: when HHI is low, genuine
bidders exert downward pressure on prices, and cover bidders serve to
create noise, satisfy minimum-bidder rules, and obscure the cartel's
bid pattern from forensic screens. When HHI is high, the cartel can
sustain prices through market power alone, reducing the marginal return
to manufactured competition.
\end{proof}

% ── 3.5 Structural Likelihood ────────────────────────────────────

\subsection{Structural Likelihood}
\label{sec:likelihood}

We specify a mixture model for the bid-level data. Each bid $b_{jt}$ in
tender $t$ is generated by one of two processes:

\begin{equation}
  b_{jt} \sim
  \begin{cases}
    F_{\text{genuine}}(b \mid \mathbf{x}_j, \boldsymbol{\alpha}) &
      \text{if firm } j \text{ is genuine (non-FL)} \\
    G_{\text{cover}}(b \mid b^*_t, \boldsymbol{\gamma}) &
      \text{if firm } j \text{ is FL}
  \end{cases}
  \label{eq:mixture}
\end{equation}

\noindent where $\boldsymbol{\alpha} = (\mu_g, \sigma_g^2)$ are the
parameters of the genuine bid distribution and $\boldsymbol{\gamma} =
(\delta \text{ or } \mu_c, \sigma_c)$ are the parameters of the
cover-bid distribution.

\paragraph{Genuine bids.} We model genuine bids as log-normal:
\begin{equation}
  \log b_{jt} \mid j \text{ genuine} \sim
  N\!\left(\mathbf{x}_j'\boldsymbol{\beta} + \alpha_i + \lambda_t,
  \; \sigma_g^2\right)
  \label{eq:genuine_dist}
\end{equation}
where $\alpha_i$ and $\lambda_t$ are item and year fixed effects, and
$\mathbf{x}_j$ includes firm size (porte), firm age, and CNAE sector
dummies.

\paragraph{Cover bids.} Under Regime~1:
\begin{equation}
  b_{jt} \mid j \text{ FL} \sim
  \text{Uniform}[b^*_t, \; b^*_t + \delta]
  \label{eq:cover_dist_r1}
\end{equation}
Under Regime~2:
\begin{equation}
  \log b_{jt} \mid j \text{ FL} \sim
  N\!\left(\log b^*_t + \epsilon, \; \sigma_c^2\right),
  \quad \sigma_c < \sigma_g
  \label{eq:cover_dist_r2}
\end{equation}

\paragraph{Likelihood.} The log-likelihood contribution of tender $t$
with $n_t$ genuine bids and $m_t$ FL bids is:
\begin{equation}
  \ell_t = \sum_{j \in \mathcal{G}_t} \log f_{\text{genuine}}(b_{jt}
  \mid \mathbf{x}_j, \boldsymbol{\alpha})
  + \sum_{\ell \in \mathcal{F}_t} \log g_{\text{cover}}(b_{\ell t}
  \mid b^*_t, \boldsymbol{\gamma})
  \label{eq:loglik}
\end{equation}
where $\mathcal{G}_t$ and $\mathcal{F}_t$ are the sets of genuine and
FL bidders in tender $t$, respectively. The FL classification provides
the mixture labels, so no EM algorithm is required---the model is a
supervised mixture with known group membership.

\paragraph{Robustness to classification error.} % V7: addresses referee concern
The FL classification is a statistical rule (median + 1.5 $\times$ IQR)
that necessarily involves misclassification: some genuine firms are
labeled FL (false positives) and some cover bidders are not (false
negatives). In a supervised mixture with noisy labels, misclassification
attenuates the estimated distributional difference between groups. If
a fraction $\alpha$ of FL-labeled bids are actually genuine, the
estimated cover-bid mean $\hat{\epsilon}$ is biased toward zero:
$\hat{\epsilon}_{\text{noisy}} = (1-\alpha) \epsilon$. The attenuation
is analogous to measurement-error bias in regression and implies that
our estimates of $\hat{\epsilon}$ and $\hat{\sigma}_c$ are conservative
(biased toward the genuine distribution). The cross-fitting exercise
(Section~\ref{sec:robustness}), which defines FL on a temporal subsample
and estimates on the complement, partially addresses this concern by
breaking the mechanical link between classification and estimation.

\paragraph{Estimation.} We estimate $(\boldsymbol{\alpha},
\boldsymbol{\gamma})$ by maximum likelihood using the bid-level data
(approximately 40 million bids total; 23.2 million non-FL losing bids
and 189,381 FL losing bids after restricting to observations with valid
prices and item-year fixed-effect coverage). Standard errors are
computed via parametric bootstrap (500 replications, parallelized
across 16 CPU cores). The genuine-bid parameters
$(\boldsymbol{\beta}, \sigma_g^2)$ are estimated from the non-FL
losing bids; the cover-bid parameters
$(\delta \text{ or } \sigma_c)$ are estimated from the FL losing
bids, conditional on the winning bid $b^*_t$ in each tender.

\paragraph{Model selection.} We select between Regime~1 and Regime~2
using the Bayesian Information Criterion (BIC) and a Kolmogorov--Smirnov
goodness-of-fit test comparing the fitted cover-bid distribution to the
empirical FL bid distribution. The regime with lower BIC and smaller KS
distance is selected as the preferred specification.

% ── 3.6 Identification ───────────────────────────────────────────

\subsection{Identification of Structural Parameters}
\label{sec:identification}

Each structural parameter is identified by specific variation in the
data:

\begin{itemize}
  \item $\boldsymbol{\beta}$, $\sigma_g^2$ (genuine bid distribution):
    identified from non-FL losing bids within item--year cells.
    The item fixed effects absorb product-specific cost levels; the year
    fixed effects absorb aggregate trends. Firm-level covariates
    $\mathbf{x}_j$ provide cross-sectional variation in costs.

  \item $\delta$ or $\sigma_c$ (cover bid parameters): identified from
    the distribution of FL bids \emph{relative to the winning bid} in
    each tender. The model requires that FL bids are measured jointly
    with $b^*_t$, which we observe directly as the negotiated price.

  \item $\theta_k$ (detection probability): partially identified from
    the cross-section of PBU characteristics. Following
    Proposition~\ref{prop:optimal_m}, $m^*$ is decreasing in
    $\theta_k$. We proxy $\theta_k$ with PBU size (number of annual
    procurements), exploiting the oversight heterogeneity
    documented in Section~\ref{sec:robustness}.
\end{itemize}

\noindent The cartel markup is not independently identified by the
distributional parameters $(\hat{\epsilon}, \hat{\sigma}_c)$---these
characterize how cover bids relate to the winning price, not the
price-setting decision itself. The markup is identified from the
reduced-form comparison of FL-present and FL-absent tenders, disciplined
by the model's prediction that FL presence should be associated with
higher prices (Prediction~\ref{pred:price_v6}) through strategic
selection into competitive markets (Prediction~\ref{pred:market_v6}).
See Section~\ref{sec:results_structural} for the consistency check
between structural and reduced-form estimates.

% ── 3.7 Testable Predictions ─────────────────────────────────────

\subsection{Testable Predictions}
\label{sec:predictions}

The model generates five testable predictions, each mapped to a specific
empirical test in Section~\ref{sec:empirical}:

\begin{prediction}[Price effect]
\label{pred:price_v6}
FL-present tenders have higher winning prices than comparable tenders
without FL, controlling for item type, year, and PBU.
\emph{Test:} OLS regression (Section~\ref{sec:reduced_form}).
\end{prediction}

\begin{prediction}[Strategic selection]
\label{pred:selection_v6}
FL-present tenders have at least as many genuine (non-FL) competitors as
FL-absent tenders. Cover bidders are deployed on top of genuine
competition, not as substitutes.
\emph{Test:} Non-FL firm count regression (Section~\ref{sec:reduced_form}).
\end{prediction}

\begin{prediction}[Regime identification]
\label{pred:regime_v6}
Under Regime~1, FL bid dispersion exceeds genuine bid dispersion
($\sigma_{\text{cover}} > \sigma_{\text{genuine}}$). Under Regime~2,
FL bid dispersion is lower ($\sigma_{\text{cover}} <
\sigma_{\text{genuine}}$). BIC model selection determines the
empirically dominant regime.
\emph{Test:} Structural likelihood comparison and bid-dispersion
regression (Section~\ref{sec:structural_estimation}).
\end{prediction}

\begin{prediction}[Market selection]
\label{pred:market_v6}
The FL--price effect is decreasing in market concentration (HHI).
Cover bidding concentrates in competitive markets where the marginal
return to manufactured competition is highest.
\emph{Test:} Network-split regression and FL $\times$ HHI interaction
(Section~\ref{sec:empirical_network}).
\end{prediction}

\begin{prediction}[Exchangeability violation]
\label{pred:exchange_v6}
FL bid residuals are non-exchangeable with non-FL bid residuals
conditional on observables. The structural model implies that FL bids
are drawn from $G_{\text{cover}}$ rather than $F_{\text{genuine}}$,
producing distinguishable residual distributions.
\emph{Test:} Kolmogorov--Smirnov test on Bajari--Ye first-stage
residuals (Section~\ref{sec:bajari_ye}).
\end{prediction}
